\section{Mathematical Formulation and Proof of Proposition 1}
\label{sec:math_proof}
The complete Mixed Integer Linear Program (MILP) for HFVRP-VC is defined as follows:
\begin{align}
\min \quad & \mathcal{F}(\mathbf{x}) \\
\text{s.t.} \quad & \sum_{k \in \mathcal{K}} \sum_{j \in \mathcal{V}_c \cup \{0\}, j \ne i} x_{ijk} = 1, \quad \forall i \in \mathcal{V}_c \label{eq:visit_once} \\
& \sum_{j \in \mathcal{V}_c} x_{0jk} = y_k, \quad \forall k \in \mathcal{K} \label{eq:depot_leave} \\
& \sum_{i \in \mathcal{V}} x_{ipk} - \sum_{j \in \mathcal{V}} x_{pjk} = 0, \quad \forall p \in \mathcal{V}_c, \forall k \in \mathcal{K} \label{eq:flow_balance} \\
& \sum_{i \in \mathcal{V}_c} q_i \sum_{j \in \mathcal{V}} x_{ijk} \le Q_k y_k, \quad \forall k \in \mathcal{K} \label{eq:capacity} \\
& \sum_{(i,j) \in \mathcal{E}} d_{ij} x_{ijk} \le R_k y_k, \quad \forall k \in \mathcal{K} \label{eq:range} \\
& u_{ik} - u_{jk} + N x_{ijk} \le N - 1, \quad \forall i, j \in \mathcal{V}_c, i \ne j, \forall k \label{eq:subtour} \\
& x_{ijk} \in \{0, 1\}, \quad y_k \in \{0, 1\}, \quad u_{ik} \ge 0
\end{align}
Constraint~\eqref{eq:visit_once} ensures every customer is visited once. Constraints~\eqref{eq:depot_leave}--\eqref{eq:flow_balance} enforce route continuity. Constraint~\eqref{eq:capacity} limits payload to vehicle capacity $Q_k$. Constraint~\eqref{eq:range} bounds distance by usable condition-aware range $R_k$. Constraint~\eqref{eq:subtour} eliminates Miller-Tucker-Zemlin (MTZ) sub-tours.

\subsection{Theoretical Guarantee: Proposition 1}
A central operational requirement in industrial logistics is that hybrid or quantum heuristic engines must never deliver a solution worse than established classical baselines.

\begin{proposition}[Conditional Non-Degradation Guardrail]
\label{prop:guardrail}
Let $S_{\text{ortools}}$ be the initial feasible solution generated by Google OR-Tools. Let $S_{\text{final}}$ be the output produced by the ROUTE/IQ-NextGen engine. Under identical road network distance matrix $\mathbf{D}$ and fleet constraints, the engine strictly satisfies:
\begin{equation}
\mathcal{D}(S_{\text{final}}) \le \mathcal{D}(S_{\text{ortools}}) + \epsilon_{\text{float}}
\end{equation}
where $\epsilon_{\text{float}} \le 10^{-4}\text{ km}$ accounts for numerical floating-point precision.
\end{proposition}

\begin{proof}
Let $\mathcal{P}_T$ be the candidate blackboard after $T$ iterations. In Stage 5 of the engine lifecycle, the centralized Guardrail evaluates the candidate with minimum distance:
\begin{equation}
S^* = \arg\min_{S \in \mathcal{P}_T \cup \{S_{\text{ortools}}\}} \mathcal{D}(S)
\end{equation}
The engine evaluates the gate condition:
\begin{equation}
\Delta = \mathcal{D}(S^*) - \mathcal{D}(S_{\text{ortools}})
\end{equation}
If $\Delta \le \epsilon_{\text{float}}$, the candidate $S^*$ is returned as $S_{\text{final}}$. If $\Delta > \epsilon_{\text{float}}$, the engine executes a deterministic fail-safe fallback:
\begin{equation}
S_{\text{final}} \leftarrow \text{copy.deepcopy}(S_{\text{ortools}})
\end{equation}
By case exhaustion, either $\mathcal{D}(S_{\text{final}}) = \mathcal{D}(S^*) \le \mathcal{D}(S_{\text{ortools}}) + \epsilon_{\text{float}}$ or $\mathcal{D}(S_{\text{final}}) = \mathcal{D}(S_{\text{ortools}})$. Therefore, $\mathcal{D}(S_{\text{final}}) \le \mathcal{D}(S_{\text{ortools}}) + \epsilon_{\text{float}}$ holds unconditionally across all executions.
\end{proof}
